\begin{abstract}
本课程项目围绕“Grover搜索算法中振幅放大过程的步进式3D展示”展开，搭建了一个支持2到3个量子比特的Grover算法玩具模型。项目重点不是求解大规模搜索问题，而是把抽象的量子态演化拆解为可观察的中间步骤：Hadamard门制备均匀叠加态，Oracle对目标态执行相位翻转，Diffusion算子执行关于平均概率幅的反演，并在每一步后保存statevector。实现上，项目使用Qiskit构造量子电路和提取中间态，使用Streamlit构建交互界面，使用Plotly绘制概率幅柱状图、概率柱状图、目标态概率曲线以及可旋转缩放的3D步进概率历史图。实验结果表明，在3量子比特、目标态为\(|101\rangle\)的例子中，目标态概率由初始的\(12.5\%\)经一次完整Grover迭代提升至\(78.125\%\)，两次迭代后提升至\(94.53125\%\)，直观体现了Grover算法的振幅放大机制。
\end{abstract}

\textbf{关键词：} Grover搜索；振幅放大；Oracle；均值倒置；Qiskit；Streamlit；Plotly；3D可视化
